\documentclass[a4paper]{article}
\usepackage{graphicx}
\usepackage{amsmath,amssymb}
\usepackage{minted}
\usepackage{multirow}
\usepackage{float}
\usepackage{algorithm}
\usepackage{algpseudocode}
\usepackage{todonotes}
\usepackage{forest}
\usepackage{xstring}
\usepackage{xcolor}
\usepackage[most]{tcolorbox}
\usepackage{xparse}
\usepackage{natbib}
\usetikzlibrary{graphs,graphdrawing}
\usegdlibrary{layered}
% ND Mind Mirror default input files.
\input{/home/donkarlo/Dropbox/repo/nd_language_ai_project/src/nd_language_ai/formal/computer/programming/tex/latex/code_clips/external_graphviz.tex}
\input{/home/donkarlo/Dropbox/repo/nd_language_ai_project/src/nd_language_ai/formal/computer/programming/tex/latex/part/control_sequence/kind/semtag/semtag.tex}
\input{/home/donkarlo/Dropbox/repo/nd_language_ai_project/src/nd_language_ai/formal/computer/programming/tex/latex/code_clips/seehere.tex}

\title{Moving Average for Trend Detection}
\date{}
\author{Mohammad Rahmani}
\begin{document}
\maketitle

\section{Definition}
A moving average (MA) smooths a sequence by replacing each observation with the average of observations in a window. It reveals slowly changing patterns but does not formally test the presence of a trend.

\section{Formula and Methodology}
For a regularly sampled sequence $y_t$ and an integer window size $k$, the trailing simple moving average is
\[
 m_t = \frac{1}{k}\sum_{i=0}^{k-1} y_{t-i},\qquad t\ge k.
\]
A trailing window uses only current and past observations. A centered window is useful for retrospective plots but uses future observations; it must not be confused with a real-time feature.

\section{Example: Online Orders}
Suppose seven consecutive daily order counts are
\[
 82,\ 94,\ 91,\ 105,\ 107,\ 99,\ 122.
\]
Their trailing seven-day average is
\[
 m_7 = \frac{82+94+91+105+107+99+122}{7}=100.
\]
Thus the seven-day local level is 100 orders/day, not necessarily the prediction for tomorrow.

\section{SQL Example}
For a table with exactly one row per calendar day:
\begin{verbatim}
SELECT day, orders,
       AVG(orders) OVER (
         ORDER BY day
         ROWS BETWEEN 6 PRECEDING AND CURRENT ROW
       ) AS ma_7
FROM daily_sales
ORDER BY day;
\end{verbatim}
This uses seven \emph{rows}; missing dates must be treated explicitly. For the first six dates, SQL returns a shorter partial average unless incomplete windows are filtered out.

\section{Advantages}
\begin{itemize}
 \item Transparent, fast, and easy to compute.
 \item Reduces short-term fluctuations.
 \item Useful as a dashboard baseline.
\end{itemize}
\section{Limitations}
\begin{itemize}
 \item Lags real changes, especially with large windows.
 \item Can conceal local spikes and turning points.
 \item Sensitive to outliers and the choice of $k$.
 \item Does not provide statistical significance or uncertainty by itself.
\end{itemize}

\section{Practical Guidance}
Compare window lengths appropriate to the sampling frequency (e.g., seven and twenty-eight days for daily retail data). Always show the raw series alongside the smoothed estimate.

\bibliographystyle{plainnat}
\bibliography{/home/donkarlo/Dropbox/repo/nd_shared_research/bibliography/bibliography.bib}
\end{document}
